%% kapa: exam=1 round=21 year=2010 period=1 form=A subject=economics subject_ko=경제원론 seq=16 number=56
%% source: Q-Net 감정평가사 기출문제 (cst00309) artlSeq=5024181 file=extracted/1008/5024181/2010년 제21회 감정평가사 1차시험 기출문제/1교시.hwp
%% conversion: kordoc 4.24.51 -> tools/gen1.py (auto); review=proofread-sonnet; flags=none
\begin{question}{56}
다음의 IS-LM 모형에 관한 설명으로 옳지 않은 것은? (단, Y는 국민소득, C는 소비, I는 투자, G는 정부지출, L은 화폐수요, M은 화폐공급, r은 이자율이다.)
\begin{framedtext}
Y = C + I + G (생산물시장의 균형식) \par
L = M (화폐시장의 균형식) \par
C = a$_{1}$ + a$_{2}$Y (0 < a$_{2}$ < 1) \par
I = b$_{1}$ + b$_{2}$r (b$_{2}$ < 0) \par
L = d$_{1}$ + d$_{2}$Y + d$_{3}$r (d$_{2}$ > 0, d$_{3}$ < 0) \par
G와 M은 정책변수 \par
a$_{1}$, a$_{2}$, b$_{1}$, b$_{2}$, d$_{1}$, d$_{2}$, d$_{3}$ 는 모두 상수 \par
\end{framedtext}
\begin{choices}
  \item a$_{2}$의 절대값이 증가하면 정부지출이 국민소득에 미치는 영향이 커진다.
  \item b$_{2}$의 절대값이 증가하면 정부지출이 국민소득에 미치는 영향이 작아진다.
  \item d$_{2}$의 절대값이 증가하면 정부지출이 국민소득에 미치는 영향이 작아진다.
  \item d$_{3}$의 절대값이 증가하면 정부지출이 국민소득에 미치는 영향이 작아진다.
  \item d$_{3}$의 절대값이 증가하면 정부지출이 이자율에 미치는 영향이 작아진다.
\end{choices}
\end{question}
\answer{4}
